\section{Introduction}

The success of machine learning led to its widespread adoption in many aspects of our daily lives. Automatic prediction and forecasting methods are now used to approve mortgage applications or estimate the likelihood of recidivism~\citep{chouldechova2017fair}. It is thus crucial to design machine learning algorithms that are {\em fair}, i.e., do not suffer from bias against or towards particular population groups. An extensive amount of research over the last few  years has focused on two key questions: how to formalize the notion of fairness in the context of common machine learning tasks, and how to design efficient algorithms that conform to those formalizations. See e.g., the survey by \cite{chouldechova2018frontiers} for an overview.

In this paper we focus on the second aspect. Specifically, we consider the problem of {\em fair clustering} and propose efficient algorithms for solving this problem. Fair clustering, introduced in~\citep{chierichetti2017fair}, generalizes the standard notion of clustering by imposing a constraint that all clusters must be {\em balanced} with respect to specific sensitive attributes, such as gender or religion. In the simplest formulation, each input point is augmented  with one of two colors (say, red and blue), and the goal is to cluster the data while ensuring that, in each cluster, the fraction of points with the less frequent color is bounded from below by some parameter strictly greater than $0$. Chierichetti \etal proposed polynomial time approximation algorithms for fair variants of classic clustering methods, such as $k$-center (minimize the {\em maximum} distance between points and their cluster centers) and $k$-median (minimize the {\em average} distance between points and their cluster centers). To this end, they introduced the notion of {\em fairlet decomposition}: a partitioning of the input pointset into small subsets, called {\em fairlets}, such that a good balanced clustering can be obtained by merging fairlets into clusters. Unfortunately, their algorithm for computing a fairlet decomposition has running time that is at least {\em quadratic} in the number of the input points. As a result, the algorithm is applicable only to relatively small data sets.

In this paper we address this drawback and propose an algorithm for computing fairlet decompositions with running time that  is {\em near-linear} in the data size. We focus on the $k$-median formulation, as $k$-center clustering is known to be sensitive to outliers. Our algorithms apply to the typical case where the set of input points lie in a $d$-dimensional space, and the distance is induced by the Euclidean norm.\footnote{E.g., all data sets used to evaluate the algorithms in~\citep{chierichetti2017fair} fall into this category.}

To state the result formally, we need to introduce some notation. Consider a collection of $n$ points $P \subseteq \mathbb{R}^d$, where each point $p \in P$ is colored either {\em red} or {\em blue}. For a subset of points $S\subseteq P$, the {\em balance} of $S$ is defined as $\bal(S) := \min\set{{\card{S_r} \over \card{S_b}}, {\card{S_b} \over \card{S_r}}}$ where $S_r$ and $S_b$ respectively denote the subset of red and blue points in $S$.
Assuming $b<r$, a clustering $\sC=\{C_1 \ldots C_k\}$  of $P$ is {\em ($r,b$)-fair} if for every cluster $C\in \sC$, $\bal(C) \geq {b\over r}$. 
In $k$-median clustering, the goal is to find $k$ centers and partition the pointset $P$ into $k$ clusters centered at the selected centers such that the sum of the distances from each point $p\in P$ point to its assigned center (i.e., the center of the cluster to which $p$ belongs) is minimized. In the $(r,b)$-fair $k$-median problem, all clusters are required to have balance at least ${b\over r}$.
Our main result is summarized in the following theorem. 
\begin{theorem}\label{thm:main-fair-k-median}
Let $T(n,d,k)$ be the running time of an $\alpha$-approximation algorithms for the $k$-median problem over $n$ points in $\mathbb{R}^d$. 
Then there exists an $O(d\cdot n \cdot \log n + T(n,d,k))$-time algorithm that given a point set $P\subseteq \mathbb{R}^d$ and balance parameters $(r,b)$, computes a $(r,b)$-fair $k$-median of $P$ whose cost is within a factor of $O_{r,b}(d \cdot \log n + \alpha)$  from the optimal  cost of $(r,b)$-fair $k$-median of $P$.
% where $O_{r,b}
\end{theorem}

The running time can be reduced further by applying dimensionality reduction techniques, see, e.g., \citep{cohen2015dimensionality, makarychev2018performance} and the references therein.

We complement our theoretical analysis with empirical evaluation. Our experiments show that the quality of the clustering obtained by our algorithm is comparable to that of~\cite{chierichetti2017fair}. At the same time, the empirical runtime of our algorithm scales almost linearly in the number of points, making it applicable to massive data sets (see Figure~\ref{fig:runtime}).  

\iffalse
By applying Johnson–Lindenstrauss transform and in particular the recent result of~\cite{makarychev2018performance} on the performance of such transforms for $k$-median clustering:
\begin{corollary}\label{cor:high-dimension-fair-k-median}
There exists an $O(n \cdot \log k \cdot \log n + t_{\textrm{k-median}})$ time algorithm that given a point set $P$ and balance parameters $(b,r)$, computes a $(r,b)$-fair $k$-median of $P$ whose cost is within $O((r^8+b^8)\cdot \log n \cdot \log k \cdot \alpha_{\textrm{k-median}})$ factor of the optimal $(r,b)$-fair $k$-median of $P$ in expectation. 
\end{corollary}
\fi

\paragraph{Related work.} Since the original paper of %Chierichetti~\etal 
\cite{chierichetti2017fair}, there has been several followup works studying fair clustering. In  particular, \cite{rosner18privacy} and \cite{bercea2018cost} studied the fair variant of $k$-center clustering (as opposed to $k$-median in our case). Furthermore, the latter paper presented a ``bi-criteria'' approximation algorithm for $k$-median and $k$-means under a somewhat different notion of fairness. However, their solution relies on a  linear program that is a relaxation of an integer linear program with at least $n^2$ variables, one for every pair of points. Thus, their algorithm does not scale well with the input size.
Another algorithm proposed in \cite{fairclustering19}, requires solving a linear program with $nk$ variables. Due to the special structure of the LP it is plausible that it can be solved efficiently, but we are not aware of any empirical evaluation of this approach.

% {\color{red}Similarly, another algorithm proposed in \cite{fairclustering19},  requires solving $n$ linear programs with $\Omega(nk)$ variables each.}

The work most relevant to our paper is a recent manuscript by~\cite{schmidt2018fair}, which proposed efficient streaming algorithms for fair $k$-means (which is similar to  $k$-median studied here). Specifically, they give a near-linear time streaming algorithm for computing a {\em core-set}: a small subset $S \subseteq P$ such that solving fair clustering over $S$ yields an approximate solution for the original point-set $P$. 
In order to compute the final clustering, however, they still need to apply a fair clustering algorithm to the core-set. Thus, our approach is complementary to the core-set approach, and the two can be combined to yield algorithms which are both fast and space-efficient\footnote{We note, however, that since core-sets typically require assigning weights to data points, such combination requires extending the clustering algorithm to weighted pointsets. In this paper we do not consider the weighted case.}.
%{\color{red}In order to compute the final clustering, however, they resort to a variant of the algorithm of \cite{chierichetti2017fair}, or another algorithm with a running time exponential in $k$. Thus, our approach is complementary to theirs. In particular, our algorithm could be applied to the core-set generated by their method in order to reduce the space usage.} 

We note that the above algorithms guarantee constant approximation factors, as opposed to the  logarithmic factor in our paper. As we show in the experimental section, this does not seem to affect the empirical quality of solutions produced by our algorithm. Still, designing a constant factor algorithm with a near-linear running time is an interesting open problem.

\paragraph{Possible settings of $(r,b)$.}
\cite{chierichetti2017fair} gave $(r,b)$-fairlet decomposition algorithms only for $b=1$. This does not allow for computing a full decomposition of the pointset into well-balanced fairlets if the numbers of red and blue points are close but not equal (for instance, if their ratio is 9:10). One way to address this could be to downsample the larger set in order to make them have the same cardinality and then compute a $(1,1)$-fairlet decomposition. The advantage of our approach is that we do not disregard any, even random, part of the input. This may potentially lead to much better solutions, partially by allowing that the clusters are not ideally balanced. The general settings of $r$ and $b$ are also considered by~\cite{bercea2018cost, fairclustering19}.

 
\paragraph{Our techniques.} Our main contribution is to design a nearly-linear time algorithm for $(r,b)$-fairlet decomposition for any integer values of $r,b$. 
Our algorithm has two steps. First, it embeds the input points into a  tree metric called {\em HST} (intuitively, this is done by computing a quadtree decomposition of the point set, and then using the distances in the quadtree).  In the second step it  solves the fairlet decomposition problem with respect to the new distance function induced by HST. The distortion of the  embedding into the HST accounts for the $\log n$ factor in the approximation guarantee. 

Once we have the HST representation of the pointset, the high-level goal is to construct ``local'' $(r,b)$-fairlets with respect to the tree. To this end, the algorithm scans the tree in a top-down order. In each node $v$ of the tree, it greedily partitions the points into fairlets so that the number of fairlets whose points belong to subtrees rooted at different children of $v$ is minimized. In particular, we prove that minimizing the number of such fairlets (which we refer to as the \emph{Minimum Heavy Point} problem) leads to an $O(1)$-approximate $(r,b)$-fairlet decomposition with respect to the distance over the tree.  
 